% ==============================================================================
%  Entornos (cajas). Estilo sobrio, inspirado en el repositorio de referencia:
%  fondo tenue, marco fino y título en negrita dentro de la caja.
%
%  Uso:
%    \begin{definicion}[Título][etiqueta] ... \end{definicion}  (ambos opcionales)
%    \begin{teorema}[...][...]      ... \end{teorema}      (también: proposicion)
%    \begin{ejemplo}[...][...]      ... \end{ejemplo}
%    \begin{demostracion}[...]      ... \end{demostracion}
%    \begin{nota}[...]              ... \end{nota}
%    \begin{resumen}[Título]        ... \end{resumen}
% ==============================================================================

% --- Colores de las cajas --------------------------------------------------------
\definecolor{fondoDef}{RGB}{237,241,247}   % azul grisáceo muy claro
\definecolor{fondoTeo}{RGB}{236,236,238}   % gris claro
\colorlet{fondoDem}{Green!15}              % verde claro (como el repositorio)
\definecolor{fondoEje}{RGB}{246,245,241}   % gris cálido muy claro
\definecolor{fondoNot}{RGB}{249,249,249}   % casi blanco
\definecolor{marco}{RGB}{140,140,145}      % marco gris

% Colores usados en figuras
\colorlet{colRes}{Maroon}

% --- Estilo base ---------------------------------------------------------------
\tcbset{
  baseapuntes/.style={
    enhanced, breakable,
    colframe=marco, boxrule=0.4pt, arc=1pt,
    left=6pt, right=6pt, top=5pt, bottom=5pt,
    fonttitle=\bfseries, coltitle=black,
    detach title, before upper={\tcbtitle.\ },
    before skip=0.8em, after skip=0.8em,
  }
}

% --- Entornos numerados (numeración por tema: 1.1, 1.2, ...; contador común) ------
% Sintaxis: \begin{definicion}[Título opcional][etiqueta opcional]
%   p.ej.   \begin{definicion}[Correlación de Pearson][def:pearson]
%   y luego \cref{def:pearson}  ->  "definición 1.1"
\newcommand{\titulocaja}[2]{#1~\thetcbcounter\ifx\relax#2\relax\else\ (#2)\fi}
\tcbset{etiqueta/.code={\ifx\relax#1\relax\else\pgfkeysalso{label={#1}}\fi}}

\NewTColorBox[auto counter, number within=chapter, crefname={definición}{definiciones}]
  {definicion}{ O{} O{} }{baseapuntes, colback=fondoDef,
  title={\titulocaja{Definición}{#1}}, etiqueta={#2}}

\NewTColorBox[use counter from=definicion, crefname={teorema}{teoremas}]
  {teorema}{ O{} O{} }{baseapuntes, colback=fondoTeo,
  title={\titulocaja{Teorema}{#1}}, etiqueta={#2}}

\NewTColorBox[use counter from=definicion, crefname={proposición}{proposiciones}]
  {proposicion}{ O{} O{} }{baseapuntes, colback=fondoTeo,
  title={\titulocaja{Proposición}{#1}}, etiqueta={#2}}

\NewTColorBox[use counter from=definicion, crefname={ejemplo}{ejemplos}]
  {ejemplo}{ O{} O{} }{baseapuntes, colback=fondoEje,
  title={\titulocaja{Ejemplo}{#1}}, etiqueta={#2}}

% --- Entornos sin numerar -------------------------------------------------------
\newtcolorbox{demostracion}[1][]{
  baseapuntes, colback=fondoDem,
  fontupper=\setlength{\jot}{6pt},   % más separación entre líneas de align* (pasos en vertical)
  title={Demostración\ifx\relax#1\relax\else\ (#1)\fi},
  after upper={\par\hfill$\blacksquare$},
}

\newtcolorbox{nota}[1][]{
  baseapuntes, colback=fondoNot,
  title={Nota\ifx\relax#1\relax\else\ (#1)\fi},
}

\newtcolorbox{resumen}[1][Resumen]{
  baseapuntes, unbreakable,   % los resúmenes son cortos: nunca se parten entre páginas
  colback=white, boxrule=0.8pt,
  detach title, before upper={{\bfseries #1}\par\smallskip},
  title={#1},
}
